% Introduction to the EU – Text-Analysis Exam Phrases (one-pager)
% Compile: pdflatex exam-phrases.tex

\documentclass[10pt,a4paper]{article}
\usepackage[utf8]{inputenc}
\usepackage[T1]{fontenc}
\usepackage[margin=1.2cm]{geometry}
\usepackage{multicol}
\usepackage{enumitem}
\setlist{nosep,leftmargin=*}
\usepackage{xcolor}
\usepackage{microtype}
\pagestyle{empty}
\definecolor{eublue}{RGB}{0,51,153}

\begin{document}
\small
\begin{center}
{\large\bfseries\color{eublue} Introduction to the EU -- Text-Analysis Exam Phrases}\\[0.2em]
\footnotesize Use these patterns when analysing a short text (Reuters piece, speech, reading excerpt).
\end{center}
\vspace{0.3em}

\begin{multicols}{2}

\section*{\color{eublue}1. Master pattern}
\emph{``This text raises issues of \textbf{[X, Y, Z]} and must be read in the context of \textbf{[period/topic]}. The author/actor adopts a \textbf{[tone/position]} and contributes to the debate on \textbf{[course theme]} by \textbf{[main move]}.''}

\section*{\color{eublue}2. Rhetorical framework (Ethos, Logos, Pathos)}
\begin{itemize}
\item \textbf{Ethos} – Who speaks? With what credibility (refugee, politician, scholar, institution)?
\item \textbf{Logos} – What is the logical structure? Which causes/effects, comparisons, categories?
\item \textbf{Pathos} – What emotions, fears, hopes are mobilised? For which audience?
\end{itemize}
\emph{``In terms of rhetorical appeals, the text combines \textbf{ethos} (first-hand witness / institutional authority), \textbf{logos} (causal argument about [X]) and \textbf{pathos} (emotional language around [Y]).''}

\section*{\color{eublue}3. Context and course blocks}
\begin{itemize}
\item \textbf{Block 1 – History}: post-1945 order, integration milestones, crises.
\item \textbf{Block 2 – Institutions}: who speaks for Europe? twin executive, Council vs Commission vs Parliament.
\item \textbf{Block 3 – Policies}: euro, enlargement, identity, borders, migration.
\item \textbf{Block 4 – Foreign \& security}: multilateralism, international law, CFSP.
\end{itemize}
\emph{``This statement belongs to \textbf{Block [1–4]} because it addresses \textbf{[treaties / institutions / policies / CFSP]} and echoes debates we saw in [author/reading].''}

\section*{\color{eublue}4. International law, multilateralism, EU unity}
\emph{``The reference to \textbf{international law} invokes the post-1945 order (UN Charter, human-rights regime). By insisting on \textbf{multilateralism} and a \textbf{united EU voice}, the text defends Europe’s identity as a community of law and a global actor which must speak consistently to remain credible.''}

\section*{\color{eublue}5. Linking to readings}
\emph{``This passage resonates with \textbf{[Zweig / Judt / Krawatzek \& Pestel / Anderson / Vinen / Bijsmans / Gilbert / Kenny \& Pearce]} because both highlight \textbf{[theme]}: for example, [author] shows how \textbf{[historical point]}, while this text updates the debate in the context of \textbf{[Brexit / Euroscepticism / institutions / borders]}.''}

\section*{\color{eublue}6. ``Analyse, don’t summarise''}
\begin{itemize}
\item Avoid retelling the story; instead:
  \begin{enumerate}
  \item Identify \textbf{who} speaks and \textbf{when/where}.
  \item Extract \textbf{key concepts} (international law, sovereignty, multilateralism, unity, borders, identity).
  \item Link them to \textbf{course blocks} and \textbf{readings}.
  \item Comment on \textbf{what is at stake} (for EU order, legitimacy, security).
  \end{enumerate}
\end{itemize}

\section*{\color{eublue}7. Stock openers}
\begin{itemize}
\item \emph{``The text comments on \textbf{[event]} and raises questions about \textbf{[law / institutions / membership / foreign policy]}.''}
\item \emph{``From the point of view of \textbf{EU integration history}, this reflects the tension between \textbf{[sovereignty vs supranationalism / unity vs diversity]}.''}
\item \emph{``Seen through the lens of \textbf{Block [2/3/4]}, the passage illustrates the problem of \textbf{[who speaks for Europe / how EU acts externally / how borders are drawn]}.''}
\end{itemize}

\section*{\color{eublue}8. Closing sentence}
\emph{``In summary, the text does not just describe \textbf{[event]}; it intervenes in the debate on \textbf{[course theme]} by defending a particular vision of \textbf{Europe’s history / institutions / borders / role in the world}. This is why it matters for understanding the current EU.''}

\end{multicols}
\end{document}

